\begin{tablalafrox}{Diccionario del Producto 9: Implantación y usuarios habilitados. Fuente: elaboración propia de los paquetes de implantación y habilitación de usuarios a partir del plan del Formulario T-18.}{tab:T14-dic-prod-9}
  {L{1.2cm} Y Y Y C{0.9cm} L{2.3cm}}
  {\cab{Código} & \cab{Paquete de trabajo} & \cab{Entregable} & \cab{Criterio de aceptación} & \cab{Resp.} & \cab{Período}}
  \multicolumn{6}{>{\columncolor{lafroxBg}\RaggedRight\arraybackslash}p{\dimexpr\textwidth-2\tabcolsep\relax}}{\textbf{Cuenta 9.1 — Planes de implantación y reversión.} Esta cuenta reúne los planes de implantación y reversión. La Etapa 1 comprende las olas por sitio y el procedimiento de reversión; la Etapa 2 cuenta con su propio plan fuera de la ventana de 05:30 a 07:00.} \tabularnewline*
  9.1.1 & Plan de olas de la Etapa 1 por sitio y zona & Plan de implantación por sitio y perfil (Talca, Concepción y plataformas de cross-docking). & Cada ola tiene su sitio, su zona, su criterio de avance escrito y su registro oficial. & IMP & E1 · Antes del H6. \\
  9.1.2 & Procedimiento y ensayo de reversión de la Etapa 1 & El procedimiento escrito y ensayado utiliza la versión operativa local para la reversión de la Etapa 1. & Ensayo de reversión con tiempos medidos, 96 despachos sin interrupción, DTE válidos emitidos sólo por ERP, WMS sólo lectura y cero pérdida/duplicación; final antes de 05:30 y fechas permitidas. & IMP & E1 · Antes de cada corte. \\
  9.1.3 & Plan de implantación y reversión de la Etapa 2 & La implantación y la reversión de la Etapa 2 se ejecutan fuera de la ventana de 05:30 a 07:00. & La reversión desactiva la capacidad de la Etapa 2 que falla sin tocar ni degradar la Etapa 1. El ensayo se realiza fuera de la ventana de 05:30 a 07:00 y de las fechas prohibidas, con tiempos medidos y sin pérdida ni duplicación de datos. & IMP & E2 · Antes del H11. \\
  \midrule
  \multicolumn{6}{>{\columncolor{lafroxBg}\RaggedRight\arraybackslash}p{\dimexpr\textwidth-2\tabcolsep\relax}}{\textbf{Cuenta 9.2 — Marcha blanca.} Esta cuenta reúne la operación supervisada de las dos marchas blancas: la de la Etapa 1 comienza con H6 y la de la Etapa 2 con H11. La segunda convive con la Etapa 1 en producción, con una sola fuente de verdad para los datos compartidos.} \tabularnewline*
  9.2.1 & Marcha blanca de la Etapa 1 con conciliación diaria (H6) & Informes diarios de indicadores y de conciliación; informe mensual de restauración de respaldos; acta de cierre de la marcha blanca. & Se cumplen las seis condiciones del Art. 17.3; cada mes se restaura una muestra representativa sin errores y se mide el tiempo efectivo (BTT, cap. 7, RT-07.12, p. 18). La Contraparte Técnica firma el acta de cierre tras revisar la evidencia del mes 15, lo que habilita el paso a producción del mes 16. & IMP & E1 · Meses 13 a 15. Empieza con el H6. \\
  9.2.2 & Marcha blanca de la Etapa 2 en convivencia con la Etapa 1 y fuente única de datos (H11) & Informes diarios de indicadores y de conciliación; informe mensual de restauración de respaldos; acta de cierre de la marcha blanca. & Se cumplen las seis condiciones del Art. 17.3, no hay doble digitación y cada mes se restaura una muestra representativa sin errores con su tiempo medido (BTT, cap. 7, RT-07.12, p. 18). La Contraparte Técnica firma el acta de cierre tras revisar la evidencia del mes 20, lo que habilita el paso a producción del mes 21. & IMP & E2 · Meses 19 y 20. Empieza con el H11. \\
  \midrule
  \multicolumn{6}{>{\columncolor{lafroxBg}\RaggedRight\arraybackslash}p{\dimexpr\textwidth-2\tabcolsep\relax}}{\textbf{Cuenta 9.3 — Estabilización.} Esta cuenta reúne la estabilización posterior al paso a producción. La Etapa 1 incluye acompañamiento en ruta y restauraciones de soporte puente; la Etapa 2 contempla apoyo a Comercial, Finanzas, cadenas y transportistas.} \tabularnewline*
  9.3.1 & Estabilización de la Etapa 1, acompañamiento en ruta y restauraciones de soporte puente & Informe de estabilización por ola y tres informes mensuales de restauración, de los meses 16 a 18. & Cada ruta recibe al menos un acompañamiento, no quedan defectos críticos ni altos abiertos y cada restauración funciona sin errores con su tiempo medido (BTT, cap. 7, RT-07.12, p. 18). & IMP & E1 · Meses 16 a 20, desde el paso a producción del 17-05-2028, tras el acta de cierre de 9.2.1 que acepta el H7 (D-28). \\
  9.3.2 & Estabilización de la Etapa 2 con Comercial, Finanzas, cadenas y transportistas & Informe de estabilización. & No quedan defectos críticos ni altos abiertos al cierre del período. & IMP & E2 · Desde el paso a producción de la Etapa 2 hasta el cierre del mes 22, con cuatro semanas reforzadas desde ese paso; con inicio contractual supuesto en febrero de 2027, del 18 de octubre al 30 de noviembre de 2028, con refuerzo del 18 de octubre al 16 de noviembre. \\
  \midrule
  \multicolumn{6}{>{\columncolor{lafroxBg}\RaggedRight\arraybackslash}p{\dimexpr\textwidth-2\tabcolsep\relax}}{\textbf{Cuenta 9.4 — Aceptación y paso a producción.} Esta cuenta reúne los expedientes de aceptación H7 y H12, con la evidencia de marcha blanca y el registro de producción de cada etapa. La garantía de correcto funcionamiento integra las condiciones de aceptación final del H12.} \tabularnewline*
  9.4.1 & Expediente y acta de aceptación de la Etapa 1 (H7) & Expediente con evidencia de 9.2.1 y registro de producción. & Firma que habilita el corte y registro de producción efectiva que completa el expediente H7 en el mes 16; secuencia conforme al T-15, sección 5.5. & JP & E1 · Mes 16: entrega el primer día hábil, acta tras diez días hábiles y paso a producción el día siguiente. \\
  9.4.2 & Expediente y acta de aceptación final (H12) & Expediente con evidencia de 9.2.2 y registro de producción. & Firma que habilita el corte y registro de producción efectiva que completa el expediente H12 en el mes 21 con garantía constituida; secuencia conforme al T-15, sección 5.5. & JP & E2 · Mes 21: entrega el primer día hábil, acta tras diez días hábiles y paso a producción el día siguiente. \\
  9.4.3 & Garantía de correcto funcionamiento constituida (H12) & Garantía constituida. & Está entregada antes de la entrega del expediente del H12. & JP & E2 · Mes 20, en paralelo con la marcha blanca de la Etapa 2; es un trámite financiero y no interviene sistemas. \\
  \midrule
  \multicolumn{6}{>{\columncolor{lafroxBg}\RaggedRight\arraybackslash}p{\dimexpr\textwidth-2\tabcolsep\relax}}{\textbf{Cuenta 9.5 — Capacitación y certificación de usuarios.} Esta cuenta reúne la capacitación de cada grupo según su rol y situación, junto con las pruebas de certificación de usuarios, una de las seis condiciones del artículo 17.3 de las Bases Administrativas (p. 13).} \tabularnewline*
  9.5.1 & Plan de capacitación por perfil & Plan de capacitación aprobado por perfil (usuarios finales, avanzados, administradores, equipo técnico y soporte). & Cubre todos los perfiles de cada ola y está aprobado por el CLIENTE. & IMP & E1 · Antes del H6. \\
  9.5.2 & Programa de capacitación continua del turno de noche & Programa de inducción continua y tutores formados para atender la rotación del 38 \% del turno de noche. & Un preparador nuevo opera su flujo después de un máximo de dos horas de aprendizaje. & IMP & Ambas · Desde el H6, y durante toda la Operación. \\
  9.5.3 & Habilitación de conductores de transportistas (escaneos y protocolo de frío) & Inducción y habilitación registrada de cada conductor. & Incluye salida de cámara, carga y bajada, rojo en TC58e y prohibición de liberar; nadie sale a ruta sin inducción. & IMP & Ambas · Meses 11–20; antes de la salida de cada conductor. \\
  9.5.4 & Traspaso al equipo de TI de cuatro personas del CLIENTE & Acuerdo de qué funciones se traspasan y qué funciones quedan como servicio, y certificación de las cuatro personas. & Las cuatro personas están certificadas en sus funciones antes del cierre de cada marcha blanca. & IMP & Ambas · Antes del H7 y del H12. \\
  9.5.5 & Materiales de capacitación por perfil & Materiales por perfil. & Cada perfil tiene su manual, su guía rápida, sus preguntas frecuentes y sus videos, entregados en formato editable. & IMP & Ambas · Antes del H6, y actualizados antes del H11. \\
  9.5.6 & Plataforma de autoformación en línea con registro de avance & Plataforma habilitada con los cursos por perfil. & Cada perfil tiene su curso, y el avance de cada persona queda registrado. & IMP & Ambas · Antes del H6. \\
  9.5.7 & Certificación de preventistas, conductores y preparadores (Etapa 1) & Registro de certificaciones por persona. & El 100 \% de los usuarios de la ola está certificado antes de su cierre. & IMP & E1 · H6 y H7. \\
  9.5.8 & Certificación de Comercial y Finanzas en costo de servir (Etapa 2) & Registro de certificaciones. & El 100 \% de los usuarios definidos está certificado antes del mes 19. & IMP & E2 · Antes del H11. \\
  \midrule
  \multicolumn{6}{>{\columncolor{lafroxBg}\RaggedRight\arraybackslash}p{\dimexpr\textwidth-2\tabcolsep\relax}}{\textbf{Cuenta 9.6 — Gestión del cambio.} Esta cuenta de control reúne el acompañamiento en terreno de los perfiles que requieren apoyo adicional a la capacitación. El diagnóstico y los agentes de cambio apoyan la Etapa 1; la Etapa 2 incorpora la comunicación con cadenas y transportistas, mientras la adopción se mide por ola.} \tabularnewline*
  9.6.1 & Plan y registros de acompañamiento a personal de larga trayectoria & Registro del acompañamiento por persona (personal con 20 o 30 años en la compañía). & Cada persona identificada tiene su plan de acompañamiento y queda certificada. & IMP & Ambas · Antes del H6 y hasta el H7. \\
  9.6.2 & Plan de gestión del cambio con cadenas y transportistas & Plan de comunicación externa ejecutado. & Cada cadena y cada transportista recibe la comunicación y su acceso antes de la marcha blanca de la Etapa 2. & IMP & E2 · Antes del H11. \\
  9.6.3 & Diagnóstico del impacto organizacional por área y perfil, con agentes de cambio & Diagnóstico de impacto y nómina de agentes de cambio habilitados. & Cada área tiene su impacto y sus resistencias previstas, con una estrategia para cada una, y cada sitio tiene al menos un agente de cambio. & IMP & E1 · Meses 2 a 6. \\
  9.6.4 & Procedimientos operativos rediseñados, contingencia de frío y ruptura confirmada & Procedimientos por ruta y proceso aprobados. & Incluye reparación, trasbordo o retorno, entrega parcial, ajuste ERP, reposición, área de no aptos y disposición; procedimiento de frío antes del piloto. & IMP & Ambas · Meses 8–20; antes del piloto y de cada ola. \\
  9.6.5 & Informes de adopción y plan de acción correctiva & Informe de adopción por ola y plan de acción cuando corresponde. & Cada ola tiene su medición, y cada meta no alcanzada tiene su plan. & IMP & Ambas · Durante cada marcha blanca y en la estabilización. \\
  \midrule
  \multicolumn{6}{>{\columncolor{lafroxBg}\RaggedRight\arraybackslash}p{\dimexpr\textwidth-2\tabcolsep\relax}}{\textbf{Cuenta 9.7 — Acuerdos con terceros.} Esta cuenta reúne los acuerdos con terceros necesarios para la salida de la solución a la ruta.} \tabularnewline*
  9.7.1 & Acuerdos operacionales firmados con los diez transportistas & Diez acuerdos firmados. & Los diez acuerdos están firmados antes de la ola de reparto. & IMP & E1 · Antes del H6. \\
  9.7.2 & Acta de acuerdo con el sindicato sobre terminales, GPS y cámaras & Acta firmada por el sindicato y el CLIENTE. & Está firmada antes de la ola de reparto, y la configuración respeta lo acordado. & IMP & E1 · Antes del H6. \\
\end{tablalafrox}
